\section{EA Test 6A: dam break past a building, laboratory scale}

Test~6A of the Environment Agency benchmark \cite{NeelzPender2013} is the
IMPACT isolated-building dam break in the flume at the Universit\'e
catholique de Louvain (Soares-Fraz\~ao and Zech \cite{SoaresFrazaoZech2002}).
A channel 3.6~m wide, with side slopes 0.34~m wide and 0.155~m high along
both walls, holds 0.4~m of water in a 6.75~m reservoir behind a 0.8~m
thick dam with a 1~m opening, and 0.02~m beyond it. The gate is lifted at
$t = 0$ and the flood strikes a 0.8~m~$\times$~0.4~m building set at
63$^\circ$ to the channel (Figure~\ref{fig:ea6a_setup}). Manning's
$n = 0.01$, the specified resolution is 0.1~m (here triangles of at most
0.01~m$^2$, about 56,000 in all), and the run lasts 2~min. The geometry is
built from the benchmark's DEM exactly as for Test~6B,
every length divided by 20.

\begin{figure}[h]
\begin{center}
\includegraphics[width=0.9\textwidth]{setup.png}
\end{center}
\caption{EA Test 6A: the flume, and the reservoir, dam, building and gauges
G1--G6 close up.}
\label{fig:ea6a_setup}
\end{figure}

The report shows the UCL measurements, available for the first 30~s,
against the packages (its Figures 4.25--4.26). Every package matched the
level at G6, upstream of the gate, which follows the periodic flow through
the gate with a period of about 9~s. Elsewhere none of them reproduced the
transcritical flow consistently: ``at best there was usually a discrepancy
in the timing of transitions''. \anuga{} behaves the same way. The level
and speed at G6 and the levels at G3 and G4 follow the measurements
closely, but the hydraulic jump that forms in front of the building and
travels back upstream reaches G2 at about 30~s with the second-order
schemes, where the measurements show it at 13--18~s and most packages at
10--12~s, and the levels at G1 and in the wake (G5) are 3--4~cm low. The
first-order scheme (DE0) puts the jump at G2 at 17~s. Friction, the mesh
(down to 0.0025~m$^2$ triangles) and the representation of the building
leave the timing unchanged; flattening the side slopes moves the answer
further from the measurements.

The checks (Figure~\ref{fig:ea6a_checks}) are the volume balance, the
order in which the bore reaches the gauges, the period of the gate flow,
a bound on the largest speed, and a stored baseline per flow algorithm.

\begin{figure}
\begin{center}
\includegraphics[width=0.95\textwidth]{stage_plot.png}
\end{center}
\caption{EA Test 6A: water level at G1--G6 over the first minute.}
\end{figure}

\begin{figure}
\begin{center}
\includegraphics[width=0.95\textwidth]{speed_plot.png}
\end{center}
\caption{EA Test 6A: speed at G1--G6 over the first minute.}
\end{figure}

\begin{figure}
\begin{center}
\includegraphics[width=0.9\textwidth]{checks.png}
\end{center}
\caption{EA Test 6A: the checks and their outcome.}
\label{fig:ea6a_checks}
\end{figure}
